% ══════════════════════════════════════════════════════════════════════════════
% Cover Letter — Submission to Mediterranean Archaeology and Archaeometry (MAA)
% Author: Seth Tenenbaum
% Date:   April 2026
% ══════════════════════════════════════════════════════════════════════════════

\documentclass[12pt,a4paper]{letter}
\usepackage[margin=1in]{geometry}
\usepackage{hyperref}
\hypersetup{colorlinks=true,urlcolor=blue!60!black}

\signature{Seth Tenenbaum\\Independent Scholar\\
\href{mailto:sethtenenbaum1@gmail.com}{sethtenenbaum1@gmail.com}\\
ORCID: \href{https://orcid.org/0009-0008-5797-2498}{0009-0008-5797-2498}}

\address{Seth Tenenbaum\\
Independent Scholar\\
\href{mailto:sethtenenbaum1@gmail.com}{sethtenenbaum1@gmail.com}}

\begin{document}

\begin{letter}{The Editors\\
Mediterranean Archaeology and Archaeometry (MAA)\\
\href{mailto:contact@maajournal.com}{contact@maajournal.com}}

\opening{Dear Editors,}

I am pleased to submit the manuscript \textbf{``A Global 3\textdegree{}
Periodic Structure in the UNESCO World Heritage Longitude Distribution:
Domed Monuments, the Babylonian Beru, and a Levantine Harmonic Grid''} for
consideration as a Research Article in \textit{Mediterranean Archaeology
and Archaeometry}.

\textbf{Topical relevance to MAA.}
This study applies archaeometric and quantitative methods to a problem rooted
in Babylonian metrology and Levantine archaeoastronomy --- two areas at the
core of MAA's scope. The primary analytical unit is the \textit{beru}, a
cuneiform angular unit (MUL.APIN, c.~1000~BCE) whose 0.1-beru subdivision
(3\textdegree{}) is tested as a global spacing signal in the UNESCO World
Heritage longitude distribution. The strongest geographic signal is centred
on the 35\textdegree{}E Levantine corridor --- the Mediterranean-to-Near
East axis that defines MAA's regional remit. The paper directly engages the
ICOMOS-IAU Thematic Study on Heritage Sites of Astronomy and Archaeoastronomy
published in the context of the UNESCO World Heritage Convention.

\textbf{Summary of findings.}
The UNESCO Cultural and Mixed World Heritage corpus ($N = 1011$ sites)
exhibits a global 3\textdegree{} periodic structure in its longitude
distribution whose phase aligns with the Levantine corridor and whose
spacing matches the metrologically attested 0.1-beru unit. Two
pre-specified primary tests survive Bonferroni correction ($k = 2$):
(1)~domed and spherical monuments (stupas, tholoi, domed shrines) are
significantly enriched at harmonic longitudes (permutation $p < 0.001$;
$3.06\times$ enrichment); and (2)~A+-bearing harmonics host $3\times$ more
total UNESCO sites than non-A+-bearing harmonics (permutation $p = 0.003$).
All results are treated as exploratory; no pre-registration was filed.

\textbf{Methodological transparency.}
The study employs three independent simulation nulls, leave-one-out
stability checks, spatial-autocorrelation correction, and a comprehensive
FDR audit across all reported tests. A circularity containment protocol
(Section~1.1) explicitly defines the limits of inference arising from the
anchor-discovery sequence. The complete analysis pipeline and data are
publicly archived on GitHub and Zenodo (DOI: 10.5281/zenodo.19520502).

\textbf{Originality.}
This is the first study to test a Babylonian metrological unit as a
longitude-spacing hypothesis against the externally curated UNESCO corpus,
and the first to document morphological specificity (dome/stupa enrichment)
within a global heritage longitude distribution.

\textbf{Ethical compliance.}
This study involves no human subjects, no animal experiments, and no
cultural objects subject to UNESCO ownership regulations. The UNESCO
XML dataset is used solely for research purposes. No conflict of interest
exists. The manuscript has not been published previously and is not under
consideration elsewhere.

\textbf{Suggested reviewers.}
I suggest the following four independent reviewers with expertise in
archaeoastronomy, ancient metrology, and spatial statistics as applied to
heritage corpora. None has collaborated with me, reviewed this manuscript,
or works at the same institution.

\begin{enumerate}
  \item \textbf{[Name, Affiliation, Email]} ---
        expertise in Babylonian astronomy / cuneiform metrology.
  \item \textbf{[Name, Affiliation, Email]} ---
        expertise in archaeoastronomy / spatial statistics.
  \item \textbf{[Name, Affiliation, Email]} ---
        expertise in Mediterranean / Near Eastern archaeometry.
  \item \textbf{[Name, Affiliation, Email]} ---
        expertise in UNESCO heritage corpus analysis / quantitative archaeology.
\end{enumerate}

% ── INSTRUCTIONS TO AUTHOR ──────────────────────────────────────────────────
% BEFORE SENDING: Fill in the four reviewer slots above (name, institution, email).
% MAA guidelines: reviewers must NOT be at the same institution as you, must NOT
% have substantially collaborated with you, and must have expertise in your content
% area/method. Suggest up to 4; you may also request exclusion of 1–2 individuals.
%
% APC WAIVER REQUEST: Append the following paragraph if applying for a waiver:
%
%   "As an independent scholar conducting unfunded research, I respectfully request
%    consideration for a partial Article Processing Charge (APC) waiver. This study
%    was carried out without institutional support or external funding, and the full
%    APC would present a significant barrier to publication. I am happy to provide
%    any documentation required to support this request."
%
% ── END INSTRUCTIONS ────────────────────────────────────────────────────────

I look forward to your consideration. Please do not hesitate to contact me
with any questions.

\closing{Sincerely,}

\end{letter}
\end{document}
